\subsection{换环；Macaulay 对偶}

命题 5.—— 设 \(\rho: A \to B\) 为 Noether 局部环之间的一个局部同态，使得剩余域扩张 \(\kappa_A \to \kappa_B\) （由 \(\rho\) 导出）有有限次数。设 \(I_A\) 为一个 Matlis A-模。

\begin{enumerate}
\def\labelenumi{\alph{enumi})}
\item
  设 \(I_{B}\) 是 \(\mathrm{Hom}_{\mathrm{A}}(\mathrm{B},\mathrm{I}_{\mathrm{A}})\) 的一个子 B-模，它由满足下列条件的从 B 到 \(I_{A}\) 中的 A-同态组成：其核包含 \(m_{B}\) 的一个幂。则 \(I_{B}\) 是一个 Matlis B-模。
\item
  设 M 为一个 B-模。典范映射
\end{enumerate}

\[
\alpha : \operatorname{Hom} _ {\mathrm{B}} (\mathrm{M}, \operatorname{Hom} _ {\mathrm{A}} (\mathrm{B}, \mathrm{I} _ {\mathrm{A}})) \longrightarrow \operatorname{Hom} _ {\mathrm{A}} (\mathrm{M}, \mathrm{I} _ {\mathrm{A}})
\]

由 \(\alpha(u)(m)=u(m)(1)\) 所定义的映射诱导一个从 \(\mathrm{D}_{\mathrm{B}}(\mathrm{M})=\mathrm{Hom}_{\mathrm{B}}(\mathrm{M},\mathrm{I}_{\mathrm{B}})\) 到子 B-模 \(\mathrm{Hom}_{\mathrm{A}}^{\mathrm{cont}}(\mathrm{M},\mathrm{I}_{\mathrm{A}})\) （\(\mathrm{D}_{\mathrm{A}}(\mathrm{M})=\mathrm{Hom}_{\mathrm{A}}(\mathrm{M},\mathrm{I}_{\mathrm{A}})\) 的）上的 B-同构，该子模由满足下列条件的映射 \(f:\mathrm{M}\to\mathrm{I}_{\mathrm{A}}\) 组成：对 M 的每个元素 \(m\)，存在一个整数 \(n\geqslant0\)，使得 \(f(\mathfrak{m}_{\mathrm{B}}^{n}m)=0\)。

上面关于 \(f\) 的条件意味着 \(f\) 是连续的，只要给 \(I_A\) 赋以离散拓扑，并给 M 赋以在每个有限生成子模上诱导 \(\mathfrak{m}_B\)-进拓扑的最细拓扑；这说明了记号的合理性。同样，条件 \(g \in I_B\) 意味着 \(g\) 是连续的，只要给 \(I_A\) 赋以离散拓扑，并给 B 赋以 \(\mathfrak{m}_B\)-进拓扑。

我们来证明 a)。对每个有限长度的 B-模 M，用 T(M) 表示 B-模 \(\mathrm{Hom}_{\mathrm{A}}(\mathrm{M},\mathrm{I}_{\mathrm{A}})\)。对有限长度 B-模之间的每个 B-线性映射 \(f:M\to N\)，用 \(\mathrm{T}(f):\mathrm{T}(\mathrm{N})\to\mathrm{T}(\mathrm{M})\) 表示 B-线性映射 \(\mathrm{Hom}_{\mathrm{A}}(f,1_{\mathrm{I}_{\mathrm{A}}})\)。第 5 节的条件 FD 1) 至 FD 4) 的验证是直接的。此外，对每个有限长度的 B-模 N，我们有 \(\mathrm{long}_{\mathrm{A}}(\mathrm{N})=\mathrm{long}_{\mathrm{B}}(\mathrm{N})\left[\kappa_{\mathrm{B}}:\kappa_{\mathrm{A}}\right]\)；由于我们有 \(\mathrm{long}_{\mathrm{A}}(\mathrm{T}(\mathrm{M}))=\mathrm{long}_{\mathrm{A}}(\mathrm{M})\)（第 3 节，定理 2），我们推出 \(\mathrm{long}_{\mathrm{B}}(\mathrm{T}(\mathrm{M}))=\mathrm{long}_{\mathrm{B}}(\mathrm{M})\)，这蕴含 FD 5)。于是可以应用第 5 节的定理 3；我们有

\[
\mathrm{T} (\mathrm{B} / \mathfrak {m} _ {\mathrm{B}} ^ {n}) = \operatorname{Hom} _ {\mathrm{A}} (\mathrm{B} / \mathfrak {m} _ {\mathrm{B}} ^ {n}, \mathrm{I} _ {\mathrm{A}}),
\]

于是 Matlis B-模 \(\varinjlim\mathrm{T}(\mathrm{B}/\mathfrak{m}_{\mathrm{B}}^{n})\) 等同于子 B-模 \(\mathrm{I}_{\mathrm{B}}\) （\(\operatorname{Hom}_{\mathrm{A}}(\mathrm{B},\mathrm{I}_{\mathrm{A}})\) 的），这就证明了 a)。

我们来证明 b)。映射 \(\alpha\) 是下述典范同构的逆映射

\[
\beta : \operatorname{Hom} _ {\mathrm{A}} (\mathrm{M}, \mathrm{I} _ {\mathrm{A}}) \longrightarrow \operatorname{Hom} _ {\mathrm{B}} (\mathrm{M}, \operatorname{Hom} _ {\mathrm{A}} (\mathrm{B}, \mathrm{I} _ {\mathrm{A}}))
\]

它把 \(v \in \mathrm{Hom}_{\mathrm{A}}(\mathrm{M}, \mathrm{I}_{\mathrm{A}})\) 对应到映射 \(v'\) （从 M 到 \(\mathrm{Hom}_{\mathrm{A}}(\mathrm{B}, \mathrm{I}_{\mathrm{A}})\) ），使得 \(v'(m)(b) = v(bm)\)（A, II, 第 74 页，命题 1）。\(v'\) 取值于 \(\mathrm{I}_{\mathrm{B}}\)，必须且只须 \(v\) 属于 \(\mathrm{Hom}_{\mathrm{A}}^{cont}(\mathrm{M}, \mathrm{I}_{\mathrm{A}})\)，由此得 b)。

推论.—— a) 若 A-代数 B 是有限的，则 B-模 \(\mathrm{I_B} = \mathrm{Hom}_{\mathrm{A}}(\mathrm{B},\mathrm{I}_{\mathrm{A}})\) 是一个 Matlis B-模。

\begin{enumerate}
\def\labelenumi{\alph{enumi})}
\setcounter{enumi}{1}
\tightlist
\item
  若 B-模 M 是 Artin 的，则映射 \(\alpha\) 是一个从 \(D_{B}(M)\) 到 \(D_{A}(M)\) 上的 B-同构。
\end{enumerate}

若 A-代数 B 是有限的，则 \(\kappa_{A}\)-代数 \(B/\mathfrak{m}_{A}B\) 也是有限的，这蕴含 \(\mathfrak{m}_{A}B\) 是 B 的一个定义理想（VIII, § 3, 第 2 节，引理 2）。由于 Matlis 模 \(I_{A}\) 的每个元素都被 \(\mathfrak{m}_{A}\) 的某个幂零化，\(\mathrm{Hom}_{\mathrm{A}}(\mathrm{B}, I_{\mathrm{A}})\) 的每个元素都被 \(\mathfrak{m}_{\mathrm{B}}\) 的某个幂零化，由此得 a)。断言 b) 由 Artin 模的每个元素都被极大理想的某个幂零化这一事实得出（第 3 节，引理 4）。

命题 5 特别适用于 A 是域 \(k\) 的情形，这时可取 \(\mathrm{I}_{\mathrm{A}} = k\) ，从而 \(\mathrm{D}_{k}(\mathrm{M}) = \mathrm{Hom}_{k}^{cont}(\mathrm{M}, k)\)（「Macaulay 对偶」）。注意，假设 \([\kappa_{\mathrm{B}} : k] < +\infty\) 特别在下列情形满足：\(k\)-代数 B 是一个有限型 \(k\)-代数在某个极大理想处的局部环（A, VIII, 附录 3, 推论 1）。

更特别地，考虑一个有限型 \(k\)-代数 S，它是 N 型分次的，且使得 \(S_0\) 是一个域，即 \(k\) 的有限次扩张。可以把命题 5 应用于局部环 \(S'\) ——S 关于极大理想 \(S_+ = \bigoplus_{n > 0} S_n\) 的局部环——或者等价地，应用于它的完备化 \(\widehat{S} = \prod_{n \geqslant 0} S_n\)（III, § 1, 第 3 节，引理 2 及 § 2, 第 12 节，例 1）。此时 S-模 \(I_{\widehat{S}} = \text{Hom}_k^{cont}(\widehat{S}, k)\) 等同于

\[
\mathrm{S} ^ {* \mathrm{gr}} = \bigoplus_ {n \geqslant 0} \operatorname{Hom} _ {k} (\mathrm{S} _ {n}, k) ;
\]

对 \(s \in S\) 与 \(u \in S^{*gr}\)，元素 \(su\) ——它属于 \(S^{*gr}\) ——是内乘 \(s \lrcorner u\)（A, III, 第 156 页及第 157 页）。例如取 \(S = k[T_1, \ldots, T_d]\)，则 \(\widehat{S} = k[[T_1, \ldots, T_d]]\)。用 \((u_\alpha)_{\alpha \in \mathbf{N}^d}\) 表示 \(k\)-向量空间 \(S^{*gr}\) 的、与 S 的基 \((T^\alpha)_{\alpha \in N^d}\) 对偶的基。\(S^{*gr}\) 的 S-模结构由下列公式描述（A, III, 第 167 页）

\[
\begin{array}{l l} \mathbf {T} ^ {\beta} u _ {\alpha} = u _ {\alpha - \beta} & \text { if } \alpha \geqslant \beta , \\ \mathbf {T} ^ {\beta} u _ {\alpha} = 0 & \text { otherwise. } \end{array}
\]
% semantic-fragments: legacy-input-seam
\relax{}
